\documentclass{article}
\usepackage{amsmath, amssymb, amsthm}

\setlength{\parindent}{0pt}

\newtheorem{claim}{Claim}

\begin{document}

Suppose $n$ balls are thrown randomly into $n$ boxes, so each ball lands in each box with uniform probability. 
Also, suppose the outcome of each throw is independent of all the other throws.

\begin{claim}
    Let $X_i$ be an indicator random variable whose value is $1$ if box $i$ is empty and $0$ otherwise. The 
    closed form expression for the probability distribution of $X_i$ is 
    \[
    \Pr\{X_i = x\} =
    \begin{cases}
    \left(1 - \frac{1}{n}\right)^n, & x = 1, \\[6pt]
    1 - \left(1 - \frac{1}{n}\right)^n, & x = 0.
    \end{cases}
    \]
    and $X_1, X_2, \dots, X_n$ are dependent random variables.
\end{claim}

\begin{proof}
    Since each ball lands in each box with uniform probability, the probability that a ball lands in box $i$ is
    \[
    \Pr\{\text{ball lands in box } i\} = \frac{1}{n}
    \]
    Taking the complement,
    \[
    \Pr\{\text{ball does not land in box } i\} = 1 - \frac{1}{n}
    \]
    Since the throws are independent, the probability that box $i$ is empty is 
    \[
    \Pr\{\text{box } i \text{ is empty}\} = \left(1 - \frac{1}{n}\right)^n
    \]
    Hence
    \[
    \Pr\{X_i = 1\} = \left(1 - \frac{1}{n}\right)^n
    \]
    and 
    \[
    \Pr\{X_i = 0\} = 1 - \left(1 - \frac{1}{n}\right)^n
    \]

    \medskip

    It remains to verify whether the random variables $X_1, X_2, \dots, X_n$ are independent. If $X_1$ and $X_2$ 
    were independent, then
    \[
    \Pr\{X_1 = 1, X_2 = 1\} = \Pr\{X_1 = 1\}\Pr\{X_2 = 1\}
    \]
    We now determine whether this equality holds. Using the prior result, we have
    \[
    \Pr\{X_1 = 1\}\Pr\{X_2 = 1\} = \left(\left(1 - \frac{1}{n}\right)^n\right)^2
    \]
    To find $\Pr\{X_1 = 1, X_2 = 1\}$, we first calculate the probability that a ball 
    lands in either box $1$ or $2$.
    \[
    \Pr\{\text{ball lands in either box } 1 \text{ or } 2\} = \frac{1}{n} + \frac{1}{n} = \frac{2}{n}
    \]
    Taking the complement,
    \[
    \Pr\{\text{ball lands in neither box } 1 \text{ nor } 2\} = 1 - \frac{2}{n}
    \]
    Since the throws are independent, the probability that both boxes $1$ and $2$ are 
    empty is
    \[
    \Pr\{\text{box } 1 \text{ and } 2 \text{ are empty}\} = \left(1 - \frac{2}{n}\right)^n
    \]
    Hence
    \[
    \Pr\{X_1 = 1, X_2 = 1\} = \left(1 - \frac{2}{n}\right)^n
    \]
    However,
    \[
    \left(1 - \frac{2}{n}\right)^n \neq \left(\left(1 - \frac{1}{n}\right)^n\right)^2
    \]
    in general. For example, when $n = 2$,
    \[
    \left(1 - \frac{2}{2}\right)^2 = 0 \quad \text{while} \quad 
    \left(\left(1 - \frac{1}{2}\right)^2\right)^2 = \frac{1}{16}
    \]
    Hence the random variables $X_1, X_2, \dots, X_n$ are not independent.

    \medskip

    Hence the probability distribution of $X_i$ is
    \[
    \Pr\{X_i = x\} =
    \begin{cases}
    \left(1 - \frac{1}{n}\right)^n, & x = 1, \\[6pt]
    1 - \left(1 - \frac{1}{n}\right)^n, & x = 0.
    \end{cases}
    \]
    and $X_1, X_2, \dots, X_n$ are dependent random variables.
\end{proof}

\begin{claim}
    The probability that at least $k$ balls fall in the first box is
    \[
    \Pr\{\text{at least } k \text{ balls fall in the first box}\} \leq \binom{n}{k}\left(\frac{1}{n}\right)^k
    \]
\end{claim}

\begin{proof}
    From $n$ balls, we can choose $k$ balls by
    \[
    \binom{n}{k}
    \]
    Since each ball lands in each box with uniform probability, the probability that a ball lands in the first 
    box is
    \[
    \Pr\{\text{ball lands in the first box}\} = \frac{1}{n}
    \]
    Consider any fixed set of $k$ balls. By independence, the probability that at least $k$ balls fall in the 
    first box is
    \[
    \Pr\{\text{those } k \text{ balls land in the first box}\} = \left(\frac{1}{n}\right)^k
    \]
    Thus the probability that $k$ balls fall in the first box is
    \[
    \Pr\{k \text{ balls land in the first box}\} = \binom{n}{k}\left(\frac{1}{n}\right)^k
    \]

    \medskip

    It remains to justify why this gives an upper bound rather than an equality for the probability that at 
    least $k$ balls fall in the first box. If more than $k$ balls fall in the first box, then a single outcome 
    is counted for several different choices of $k$ balls. Hence, when summing over all $\binom{n}{k}$ choices, 
    such outcomes are counted more than once. Therefore the above expression overestimates the desired 
    probability, which gives the inequality.

    \medskip

    Hence the probability that at least $k$ balls fall in the first box is
    \[
    \Pr\{\text{at least } k \text{ balls fall in the first box}\} \leq \binom{n}{k}\left(\frac{1}{n}\right)^k
    \]
\end{proof}

\begin{claim}
    Let $R$ be the maximum of the numbers of balls that land in each of the boxes. Then 
    \[
    \Pr\{R \geq k\} \leq \frac{n}{k!}
    \]
\end{claim}

\begin{proof}
    Let $A_i$ be the event that at least $k$ balls fall into box $i$. Then
    \[
    \Pr\{R \geq k\} = \Pr\left\{\bigcup_{i=1}^{n}A_i\right\}
    \]
    By the union bound,
    \[
    \Pr\{R \geq k\} \leq \sum_{i=1}^{n}\Pr\{A_i\}
    \]
    Since each box is equally likely to receive balls,
    \[
    \sum_{i=1}^{n}\Pr\{A_i\} = n\Pr\{A_1\}
    \]

    \medskip

    From the previous claim, 
    \[
    \Pr\{A_1\} \leq \binom{n}{k}\left(\frac{1}{n}\right)^k
    \]
    Expanding the expression, 
    \[
    \begin{aligned}
    \binom{n}{k}\left(\frac{1}{n}\right)^k &= \frac{n(n - 1)\dots(n - k + 1)}{k!}\left(\frac{1}{n^k}\right) \\
    &= \frac{1}{k!}\left(\frac{n(n - 1)\dots(n - k + 1)}{n^k}\right) \\
    &= \frac{1}{k!} \cdot \prod_{j=0}^{k-1}\frac{n-j}{n}
    \end{aligned}
    \]
    Since each factor in the product is at most $1$,
    \[
    \Pr\{A_1\} \leq \frac{1}{k!}
    \]
    Therefore, 
    \[
    \Pr\{R \geq k\} \leq n\Pr\{A_1\} \leq \frac{n}{k!}
    \]

    \medskip

    Hence 
    \[
    \Pr\{R \geq k\} \leq \frac{n}{k!}
    \]
\end{proof}

\begin{claim}
    For all $\epsilon > 0$,
    \[
    \lim_{n \to \infty}\Pr\{R \geq n^\epsilon\} = 0
    \]
\end{claim}

\begin{proof}
    Substituting $k=n^\epsilon$ into the previous claim gives
    \[
    \Pr\{R \geq n^\epsilon\} \leq \frac{n}{(n^\epsilon)!}
    \]
    Since factorial growth dominates polynomial growth,
    \[
    \lim_{n \to \infty}\frac{n}{(n^\epsilon)!} = 0
    \]

    \medskip

    Hence,
    \[
    \lim_{n \to \infty}\Pr\{R \geq n^\epsilon\} = 0
    \]
\end{proof}

\end{document}
